% Source: https://artofproblemsolving.com/community/c3h1468510_venn_diagram_problem
% Stable identity: aops-topic-1468510-q01
The total number of correct answers is $90+85+60+40=275$. Let $x$ be the number of students who advanced.

\textbf{Least.} Each advancing student has at most $4$ correct answers and each other student at most $2$, so $4x+2(100-x)\ge275$, which gives $x\ge37.5$, so $x\ge38$. This is attained: let $37$ students get all four problems, $1$ student get problems $1,2,3$, and among the other $62$ students (two problems each) let $3$ get problems $1,4$, $37$ get problems $1,2$, $12$ get problems $1,3$ and $10$ get problems $2,3$. The counts per problem are $37+1+3+37+12=90$, $37+1+37+10=85$, $37+1+12+10=60$ and $37+3=40$. So the least number is $\boxed{38}$.

\textbf{Greatest.} Each advancing student has at least $3$ correct answers, so $3x\le275$ and $x\le91$. This is attained: $7$ students get nothing right, $1$ student gets only problem $1$, $1$ student gets only problem $4$, and the $91$ advancing students each miss exactly one problem: $2$ miss problem $1$, $6$ miss problem $2$, $31$ miss problem $3$ and $52$ miss problem $4$. The counts are $1+(91-2)=90$, $91-6=85$, $91-31=60$ and $1+(91-52)=40$. So the greatest number is $\boxed{91}$.
